\documentclass[11pt]{amsart}

\usepackage[margin=1.1in]{geometry}
\usepackage{amsmath,amssymb,amsthm}
\usepackage{booktabs}
\usepackage{array}
\usepackage{enumitem}
\usepackage[hidelinks]{hyperref}
\usepackage{xcolor}
\usepackage{alltt}

\newtheorem{theorem}{Theorem}[section]
\newtheorem{lemma}[theorem]{Lemma}
\newtheorem{proposition}[theorem]{Proposition}
\newtheorem{corollary}[theorem]{Corollary}
\theoremstyle{definition}
\newtheorem{definition}[theorem]{Definition}
\theoremstyle{remark}
\newtheorem{remark}[theorem]{Remark}

\emergencystretch=2em
\newcommand{\N}{\mathbb{N}}
\newcommand{\Z}{\mathbb{Z}}
\newcommand{\Mid}{\mathrm{Mid}}
\newcommand{\Ex}{\mathrm{Ex}}
\newcommand{\Sh}{\mathrm{Sh}}
\newcommand{\lean}[1]{\texttt{#1}}

\title[Erd\H{o}s Problem \#848]{Erd\H{o}s Problem \#848: an explicit sieve proof\\ and a kernel-checked Lean formalisation}
\author{macchacat}
\date{October 2026}


\begin{document}

\begin{abstract}
Let $A\subseteq\{1,\dots,N\}$ be such that $ab+1$ is not squarefree for all
$a,b\in A$ (with $a=b$ allowed). We give a proof that
$|A|\le\lfloor (N+18)/25\rfloor$ for every $N\ge1$, the number of $n\le N$ with
$n\equiv7\pmod{25}$; the bound is attained by that residue class. Sawhney proved
this for all sufficiently large $N$ without an explicit threshold, and Li gave the
first complete proof for all $N$, with a Lean formalisation. Our argument is
independent of Li's: we make every error term in Sawhney's sieve argument explicit,
strengthen the lemmas so that the analytic part covers all
$N\ge 6\times10^5$, and treat the remaining range by a colouring of the graph of
compatible integers. The whole theorem is formalised in Lean~4 with Mathlib. Every
computation is evaluated by the Lean kernel (\texttt{decide +kernel}); no
\texttt{native\_decide} is used, and the only axioms are \texttt{propext},
\texttt{Classical.choice} and \texttt{Quot.sound}. The kernel checks take about five
CPU hours and at most $9.3$~GB of memory per module.
\end{abstract}

\maketitle

\section{Introduction}

For a positive integer $N$ call a set $A\subseteq\{1,\dots,N\}$
\emph{admissible} if $ab+1$ is not squarefree for all $a,b\in A$, where $a=b$ is
allowed. Since $7\cdot7+1=50$ and $25\mid ab+1$ whenever
$a\equiv b\equiv7\pmod{25}$, the set
\[
  C_7(N)=\{n\le N:\ n\equiv7\pmod{25}\},\qquad
  c_7(N):=|C_7(N)|=\Big\lfloor\frac{N+18}{25}\Big\rfloor ,
\]
is admissible. Erd\H{o}s and S\'ark\"ozy asked whether this is extremal; this is
Problem \#848 on the Erd\H{o}s problems site~\cite{EP848}.

\begin{theorem}\label{thm:main}
For every $N\ge1$, every admissible $A\subseteq\{1,\dots,N\}$ satisfies
$|A|\le c_7(N)$. Hence the maximum size of an admissible set is exactly $c_7(N)$.
\end{theorem}

\subsection*{Prior work}
Sawhney~\cite{Sawhney} proved Theorem~\ref{thm:main} for all sufficiently large
$N$, without an explicit threshold (see also~\cite[\S IV.1]{AI2511}). The first
proof for \emph{all} $N$ is due to Li~\cite{Li,LiRepo}, who reduces the extremal
problem by a Hall-type equivalence to a completion inequality relative to the
progression $7\pmod{25}$, controls $N\le5\times10^6$ by a prefix-compatible
colouring, and formalises the complete theorem in Lean~4. The present work was
carried out independently and follows a different route: Sawhney's sieve
framework, made fully explicit.

\subsection*{This paper}
\begin{enumerate}[leftmargin=*]
\item \emph{An explicit sieve argument} (Sections~\ref{sec:notation}--\ref{sec:large}).
  Every error term in Sawhney's argument is made explicit (Lemmas~A$'$, B, C), with a
  refined case analysis in which each witness pays the large-prime error only once.
  Interval certificates in exact rational arithmetic and a tail estimate prove the
  bound for all $N\ge1.6\times10^8$.
\item \emph{Strengthened lemmas} (Section~\ref{sec:strong}). Sharper bounds for the
  large-prime set (a periodic discrepancy estimate), a Pell-type bound counted in a
  window, Bonferroni truncation, and minimal witnesses lower the analytic threshold
  to $6\times10^5$.
\item \emph{Small $N$} (Section~\ref{sec:small}). For small $N$ we exhibit a
  partition of the vertex set into $c_7(N)$ independent sets of the
  ``compatibility'' graph, built incrementally and replayed by an independent
  checker; this is done up to $N=4\times10^8$ in C\texttt{++}.
\item \emph{A kernel-checked formalisation} (Section~\ref{sec:lean}). The whole
  theorem is proved in Lean~4 with Mathlib. All computations are evaluated by the
  kernel, so neither the Lean compiler nor any external program is trusted.
  Untrusted programs only produce certificates; the kernel checks them, and the
  soundness of every checker is proved in Lean.
\end{enumerate}

\section{Notation and basic facts}\label{sec:notation}

If $a\in A$ then $a^2+1$ is not squarefree. If $p^2\mid a^2+1$ then $p\ne2$ (as
$a^2+1\not\equiv0\pmod4$) and $-1$ is a square mod $p$, so $p\equiv1\pmod4$.

Write $A_7=A\cap(7+25\Z)$, $A_{18}=A\cap(18+25\Z)$ and
$A^*=A\setminus(A_7\cup A_{18})$. A \emph{U1 prime} is a prime $p\equiv1\pmod4$
with $p\ge13$. For a U1 prime $p$ let $\pm r_p$ be the two solutions of
$x^2\equiv-1\pmod{p^2}$ and $R_p=\{x: x\equiv\pm r_p\pmod{p^2}\}$. An integer $a$
is \emph{compatible} with $b$ if $\mu(ab+1)=0$.

\begin{enumerate}[label=(F\arabic*),leftmargin=*]
\item If $x\in A^*$, then $p^2\mid x^2+1$ for some U1 prime $p$, because
  $25\mid x^2+1$ iff $x\equiv\pm7\pmod{25}$.
\item If $a\in A_7\cup A_{18}$ and $b\in A^*$, then $25\nmid ab+1$.
\item If $a\in A_7$ and $b\in A_{18}$, then $25\nmid ab+1$, since $7\cdot18+1=127$.
\item If $a,b$ are odd, then $4\mid ab+1$ iff $a\not\equiv b\pmod4$; if $a$ or $b$
  is even, $ab+1$ is odd.
\item If $p\nmid b$, then $p^2\mid ab+1$ iff $a\equiv c_p(b):=-b^{-1}\pmod{p^2}$.
  For a U1 prime $p$, $c_p(b)\equiv\pm r_p$ iff $b\equiv\pm r_p$.
\item $\rho^*(b):=\max_{(c,b)=1}\#\{s\bmod b: s^2\equiv c\}\le
  2^{\omega_{\mathrm{odd}}(b)}\cdot\{1,2,4\}$, the last factor according as
  $v_2(b)\le1$, $=2$, $\ge3$.
\end{enumerate}

\section{Sieve lemmas}\label{sec:sieve}

\subsection{Lemma A$'$: one witness, shared sets}
Fix a \emph{witness} $1\le b\le N$ and parameters $5\le z\le p_c<Y$. Define
\[
  \Mid(b)=\{a\le N:\exists p\in(p_c,Y],\ p\nmid b,\ p^2\mid ab+1\},\qquad
  \Ex(b)=\{a\le N:\exists p>Y,\ p^2\mid ab+1\}.
\]

\begin{lemma}[A$'$]\label{lem:A}
\begin{enumerate}[label=(\roman*)]
\item $|\Mid(b)|\le\sum_{p_c<p\le Y}(N/p^2+1)$.
\item $|\Ex(b)|\le\rho^*(b)\big(2(N+1)/Y+(Nb+1)/Y^2\big)$.
\item Let every prime factor $p$ of $q$ satisfy $p^2\mid q$, let $T$ be a set of
  classes mod $q$ with $p^2\nmid tb+1$ for all $t\in T$, $p\mid q$; let
  $\mathcal P=\{p\nmid qb\}$ and $k=\#\{p\in\mathcal P: p\le z\}$. With
  $S=\{a\le N: a\bmod q\in T\}$,
  \begin{multline*}
  \#\{a\in S:\ a\text{ compatible with }b\}\le
  |T|\Big[\frac Nq\Big(1-\prod_{p\in\mathcal P,\,p\le z}\Big(1-\frac1{p^2}\Big)\Big)+2^k-1\Big]\\
  +\sum_{\substack{z<p\le p_c\\ p\in\mathcal P}}|T|\Big(\frac N{qp^2}+1\Big)
  +|\Mid(b)\cap S|+|\Ex(b)\cap S|.
  \end{multline*}
\end{enumerate}
\end{lemma}

\begin{proof}
If $a$ is compatible with $b$, pick $p$ with $p^2\mid ab+1$; then $p\nmid b$, and
$p\nmid q$ since otherwise $ab+1\equiv tb+1\not\equiv0\pmod{p^2}$. Primes $p\le z$
are handled by inclusion--exclusion in each class $t$ (error $\le2^k-1$); primes
$z<p\le p_c$ by counting one class mod $qp^2$; primes $p_c<p\le Y$ give $\Mid(b)$,
and each contributes at most $N/p^2+1$ values of $a$ by (F5). For $p>Y$ write
$ab+1=p^2m$ with $m\le M=\lfloor(Nb+1)/Y^2\rfloor$; $a$ is determined by $(p,m)$,
the congruence $p^2m\equiv1\pmod b$ puts $p$ into at most $\rho^*(b)$ classes mod
$b$, and $p\le\sqrt{(Nb+1)/m}$. Summing over $m$ gives (ii).
\end{proof}

The key point is that when one witness $b$ sieves several disjoint parts of $A$,
the contributions of $\Mid(b)$ and $\Ex(b)$ add up to at most
$\Sigma(b):=|\Mid(b)|+|\Ex(b)|$: each witness pays the large-prime error once.

\subsection{Lemma B: sets of type $A^*$}
\begin{lemma}[B]\label{lem:B}
Let $q\in\{25,50,100\}$, $T$ a set of $\tau$ classes mod $q$,
$13\le z_B\le p_B\le Y'$ and $k=\#\{\text{U1 primes}\le z_B\}$. Then
$\#\{n\le N: n\bmod q\in T,\ \exists\,\text{U1 }p,\ p^2\mid n^2+1\}$ is at most
\begin{multline*}
\tau\Big[\frac Nq\Big(1-\prod_{\text{U1},\,p\le z_B}\Big(1-\frac2{p^2}\Big)\Big)+3^k-1\Big]
+\sum_{\text{U1},\,z_B<p\le p_B}2\tau\Big(\frac N{qp^2}+1\Big)\\
+\sum_{\text{U1},\,p_B<p\le Y'}2\Big(\frac N{p^2}+1\Big)
+\Big\lfloor\frac{N^2+1}{Y'^2}\Big\rfloor\big(\log_3 N+1\big).
\end{multline*}
\end{lemma}

\begin{proof}
Small primes by inclusion--exclusion over the disjoint classes $\pm r_p$; medium
primes directly. For $p>Y'$ write $n^2+1=p^2m$, so $(n,p)$ solves
$x^2-my^2=-1$. If $x_1<x_2$ are two solutions then
$x_2+y_2\sqrt m=(x_1+y_1\sqrt m)(a+c\sqrt m)$ with $a^2-mc^2=1$, $a,c\ge1$, and
$x_2\ge ax_1\ge3x_1$ unless $m=3$, which is impossible ($3\nmid n^2+1$). Hence each
$m$ contributes at most $\log_3N+1$ values of $n$.
\end{proof}

(The original form of the Pell term uses the fundamental unit, with coefficient
$1/(2\log(1+\sqrt2))$; the elementary bound above is the one formalised.)

\subsection{Lemma C: sieving $A^*$ by an even witness}
Let $b\in A^*$ be even and $T_{25}$ the 23 classes mod 25 other than 7 and 18. For
$x\in A^*$ either $x=b$ or some odd prime $p_2\nmid b$ has $p_2^2\mid xb+1$, and by
(F1) $x\in R_{p_1}$ for a U1 prime $p_1$. Writing $C_p=\{x\equiv c_p(b)\pmod{p^2}\}$,
$U_1^{\le z_1}=\bigcup_{\text{U1},\,p\le z_1}R_p$ and similarly $U_1^{>z_1}$,
\begin{equation}\label{eq:C}
|A^*|\le|T_{25}\cap U_1^{>z_1}|+\sum_{\substack{p_2\le z_2\\ p_2\text{ odd}}}
\big|T_{25}\cap U_1^{\le z_1}\cap C_{p_2}\big|
+\sum_{z_2<p_2\le p_c}\Big(\frac N{p_2^2}+1\Big)+|\Mid(b)\cap A^*|+|\Ex(b)\cap A^*|.
\end{equation}
Let $\mathcal U$ be the U1 primes $\le z_1$, $k_1=|\mathcal U|$ and
$P=\prod_{q\in\mathcal U}(1-2/q^2)$. Each term of \eqref{eq:C} is bounded
explicitly: the first by Lemma~\ref{lem:B} with $q=25$, $\tau=23$; the term
$p_2=5$ by $\frac N{25}(1-P)+3^{k_1}$ (it is one whole class mod 25, inside
$T_{25}$); an odd $p_2$ that is neither 5 nor a U1 prime by
$\frac{23}{25}\frac N{p_2^2}(1-P)+3\cdot3^{k_1}$; a U1 prime $p_2$ with
$b\in R_{p_2}$ by $\mathrm{whole}(p_2)=\frac{23}{25}\frac N{p_2^2}+3$; and a U1
prime with $b\notin R_{p_2}$ by
$\mathrm{non}(p_2)=\frac{23}{25}\frac N{p_2^2}(1-\mathrm{pr}_{p_2})+3\cdot3^{k'_{p_2}}$,
where $\mathrm{pr}_p=\prod_{q\in\mathcal U\setminus\{p\}}(1-2/q^2)$ and
$k'_p=|\mathcal U\setminus\{p\}|$.

The set $S_b=\{\text{U1 }p: b\in R_p\}$ is not known, so we bound its worst case.
Let $q_1=13<q_2=17<q_3=29<\cdots$ be the U1 primes and $g(p)=\mathrm{whole}(p)-\mathrm{non}(p)$.

\begin{lemma}[C1, C2]
$g$ is strictly decreasing on U1 primes. If
$L=\max\{\ell: q_1\cdots q_\ell\le\sqrt{N^2+1}\}$, then for every even $b\le N$
\[
\sum_{p\in S_b,\ p\le z_2}g(p)\le\sum_{j\le L,\ q_j\le z_2}\max(g(q_j),0).
\]
\end{lemma}

\begin{proof}
The monotonicity follows from the explicit form of $g$ on either side of $z_1$.
For the second claim, the primes of $S_b$ are distinct and $p^2\mid b^2+1$ for each,
so $\prod_{p\in S_b}p\le\sqrt{N^2+1}$; the $j$-th smallest element of $S_b$ is
$\ge q_j$, so $|S_b|\le L$ and $g(s_j)\le g(q_j)$.
\end{proof}

\section{Case analysis}\label{sec:cases}
Write $R(T,q;\mathrm{excl})$ for the terms of Lemma~\ref{lem:A}(iii) other than
$\Mid$ and $\Ex$, in the worst case $\mathcal P=$ all primes outside
$\mathrm{excl}$, and $B(\tau,q)$ for the bound of Lemma~\ref{lem:B}. ``Lemma C''
denotes \eqref{eq:C} without its last two terms; together with the shared part of
$A_7\cup A_{18}$ for the same witness, those terms are bounded by $\Sigma(b)$.

\begin{center}
\renewcommand{\arraystretch}{1.25}
\begin{tabular}{@{}l >{\raggedright\arraybackslash}p{4.2cm} >{\raggedright\arraybackslash}p{9.2cm}@{}}
\toprule
Case & Condition & Upper bound\\
\midrule
I.a & $A^*=\emptyset$, $A_7$ or $A_{18}$ empty & $c_7(N)$\\
I.b & $A^*=\emptyset$, $b\in A_7$, $b'\in A_{18}$ & $R(\{18\},25;\{5\})+\Sigma(b)+R(\{7\},25;\{5\})+\Sigma(b')$\\
II & $A^*$ has an even element $b$ & $R(\{7,18\},25;\{2,5\})+\Sigma(b)+\text{Lemma C}$\\
III.1 & $A^*$ odd, with $b_1\equiv1$, $b_3\equiv3\pmod4$ & $R(6\text{ cl.},100)+\Sigma(b_1)+R(2\text{ cl.},100)+\Sigma(b_3)+B(23,50)$\\
III.2.a & $A^*$ in one class mod 4, $A_7\cup A_{18}$ odd & $2(N/100+1)+R(2\text{ cl.},100)+\Sigma(b)+B(23,100)$\\
III.2.b & as III.2.a, but $A_7\cup A_{18}$ has an even $b'$ & $(N/100+1)+R(3\text{ cl.},100)+\Sigma(b)+R(4\text{ cl.},100)+\Sigma(b')+B(23,100)$\\
\bottomrule
\end{tabular}
\end{center}

The main coefficients (times 25) are $0.741$, $0.784$, $0.632$, $0.739$, $0.686$
for I.b, II, III.1, III.2.a, III.2.b; all are $<1$. The cases are exhaustive, and
the residue conditions used are consequences of (F2)--(F4). For instance in III.1,
$b_1$ sieves the classes $57,93,32,82,18,68\pmod{100}$ of $A_7\cup A_{18}$ and
$b_3$ the classes $7,43$; in each of them $4\nmid tb+1$ and $25\nmid tb+1$, so
$\{2,5\}$ may be excluded.

\section{Large $N$: interval certificates}\label{sec:large}

\begin{lemma}[D]\label{lem:D}
Fix parameters on an interval $[N_{lo},N_{hi}]$. Suppose that for every case and
every integer $N$ in the interval the bound has the form $C N+E(N)$ with
$E$ nondecreasing. If $C<1/25$ and $CN_{lo}+E(N_{hi})<(N_{lo}-6)/25$, then
$|A|<(N-6)/25\le c_7(N)$ for every integer $N\in[N_{lo},N_{hi}]$.
\end{lemma}

\begin{proof}
$(N-6)/25-CN$ is increasing, and $c_7(N)\ge(N+18-24)/25$.
\end{proof}

The parameters ($z$, $p_c$, $Y$, $Y'_q$, $p_{B,q}$, $z_B$, $z_1$, $z_2$, $L$) are
chosen by a floating-point optimiser; since every choice gives a valid bound, they
need no verification. The bounds are then recomputed in exact rational arithmetic:
$\sum1/p^2$ by integer sums of $\lceil K/p^2\rceil$ up to $2\times10^7$ and
$\sum_{n>P}1/n^2\le1/P$ beyond, $\pi(x)$ exactly up to $2\times10^7$ and by an
explicit upper bound above (in the formalised version, Chebyshev's bound
$\theta(x)\le x\log4$, which is in Mathlib), and logarithms by rational bounds in the
safe direction. In the formalised version, $2913$ intervals with integer endpoints
cover $[1.6\times10^8,10^{60}]$, with worst ratio $0.9891$ (the ratio of the
bound to $(N_{lo}-6)/25$).

For $N\ge10^{60}$ the parameters are fixed functions of $N$ (for instance
$Y=N^{0.7}$ and $p_c=p_B=2000$), $\rho^*(b)\le2N^{0.195}$ by Robin's bound on
$\omega(n)$, and $C$ is constant. Each term of $E(N)/N$ has the form
$aN^{\alpha-1}(\log N)^\gamma$ with $\alpha\le0.8$ and is decreasing for
$N\ge10^{60}$; at $N=10^{60}$ the ratios are at most $0.826$ in every case. The
formalised tail uses $t=\lfloor N^{1/4}\rfloor$, $Y=\lfloor N/t\rfloor$ and the
bound $\rho(b)^5\cdot(3\cdot5\cdots31)\le1024\cdot32^{10}\,b$, valid for all $b$,
instead of Robin's inequality. This proves Theorem~\ref{thm:main} for
$N\ge1.6\times10^8$.

\section{Strengthened lemmas: the analytic part from $6\times10^5$}\label{sec:strong}

\subsection{The large-prime set}
Fix a witness $b$ and $Y>p_c$; put $M=\lfloor(Nb+1)/Y^2\rfloor$,
$X_m=\sqrt{(Nb+1)/m}$ and $r_b(m)=\#\{s\bmod b: s^2m\equiv1\pmod b\}$.

\begin{lemma}[E]\label{lem:E}
\begin{enumerate}[label=(E\arabic*),start=0]
\item $|\Ex(b)|\le\sum_{m\le M}r_b(m)\big(\lfloor(X_m-Y)/b\rfloor+1\big)$.
\item $\sum_{m\le M}\big((X_m-Y)/b+1\big)\le(N+1)/Y+M$; hence
  $|\Ex(b)|\le\rho^*(b)\big((N+1)/Y+M\big)$.
\item Let $e=\min(v_2(b),3)$, $Q$ a set of odd primes dividing $b$, $P=2^e\prod Q$,
  $g_P=r_P$, $\alpha_P=\varphi(P)/P$ and
  $\beta_P=\max_{0\le R<P}\big(\sum_{m\le R}g_P(m)-\alpha_PR\big)$. Then
  \[
  |\Ex(b)|\le2^{\omega_{\mathrm{odd}}(b)-|Q|}\Big[\alpha_P\Big(\frac{N+1}Y+M\Big)
  +\beta_P\Big(\frac{\sqrt{Nb+1}}b+1\Big)\Big].
  \]
\end{enumerate}
\end{lemma}

\begin{proof}[Sketch]
(E0) $a\mapsto(p,m)$ with $ab+1=p^2m$ is injective and $p$ lies in $r_b(m)$ classes
mod $b$ in $(Y,X_m]$. (E1) follows from $\sum_{m\le M}m^{-1/2}\le2\sqrt M-1$ and
$(K-\sqrt M)^2\ge0$ with $K=\sqrt{Nb+1}/Y$. For (E2), $r_b\le2^{\omega_{\rm odd}(b)-|Q|}g_P$
by the Chinese remainder theorem; $\sum_{m\bmod P}g_P(m)=\varphi(P)$, so
$G(R)=\sum_{m\le R}g_P(m)\le\alpha_PR+\beta_P$; Abel summation against the
nonincreasing $f(m)=\lfloor(X_m-Y)/b\rfloor+1$ and (E1) finish the proof.
\end{proof}

To use (E2) uniformly over unknown witnesses, a witness $b\le x$ is described by
its \emph{structure} $(e,S,u)$: $e=\min(v_2(b),3)$, $S$ the small odd primes
dividing $b$, $u$ the number of larger odd prime factors. Each structure has a
least possible value $b_{\min}$, and the bound is maximised over the structures
with $b_{\min}\le x$, using the table of $(P,\varphi(P),P\beta_P)$.

\begin{lemma}[M]
A class mod $p^2$ with $p>\sqrt N$ has at most one element in $[1,N]$, so
$|\Mid(b)|\le\sum_{p_c<p\le\min(Y,\sqrt N)}(N/p^2+1)+\#\{p:\max(p_c,\sqrt N)<p\le Y\}$.
\end{lemma}

\subsection{The U1 tail}
\begin{lemma}[H]
For a U1 prime $p$, $\#\{n\le N: p^2\mid n^2+1\}\le2(N/p^2+1)$; it is $\le2$ if
$p^2>N$ and $\le1$ if $p^2>2N$.
\end{lemma}

\begin{lemma}[P]
Let $Y_q\ge5$, $M=\lfloor(N^2+1)/Y_q^2\rfloor$, and call $m$ \emph{admissible} if
$m\bmod8\in\{1,2,5\}$ and $3,7,11,19,23\nmid m$. Then
\[
\#\{n\le N:\exists\,p\ge Y_q,\ p^2\mid n^2+1\}\le\sum_{\substack{m\le M\\ m\text{ admissible}}}
\Big(1+\Big\lfloor\frac{\log(N/x_m)}{\log g_m}\Big\rfloor_+\Big),
\]
where $x_m=\sqrt{mY_q^2-1}$ and $g_m=\max(3,2\sqrt m)$.
\end{lemma}

\begin{proof}
With $n^2+1=p^2m$ we have $m\equiv n^2+1\in\{1,2,5\}\pmod8$, no prime
$\equiv3\pmod4$ divides $m$, and $n\ge x_m$. Consecutive solutions of
$x^2-my^2=-1$ grow by a factor $\ge g_m$, as in Lemma~\ref{lem:B}.
\end{proof}

\subsection{Bonferroni truncation and minimal witnesses}
\begin{lemma}[Bf]
Let $D$ be a set whose intersection with each class mod $qQ$ ($Q$ odd, coprime to $q$)
has size in $[L/(qQ)-\mu,L/(qQ)+\mu]$, let $p_1,\dots,p_k$ be primes coprime to $q$,
and let $E_i$ consist of $c$ classes mod $p_i^2$, $v_i=c/p_i^2$. For odd $d$, the
number of elements of $D$ in one class mod $q$ lying in $\bigcup E_i$ is at most
$\frac Lq\sum_{j=1}^d(-1)^{j+1}e_j(v_1,\dots,v_k)+\mu\sum_{j=1}^d\binom kjc^j$, and the
main term is nondecreasing in the set of primes.
\end{lemma}

\emph{Minimal witnesses.} Each witness is taken to be the least element of its part,
so the main terms may use $L=N-x+1$ in place of $N$. The witnesses are grouped into
buckets $[x_a,x_b]$ with $x_b=\max(x_a+1,\lfloor1.3x_a\rfloor)$, and in each bucket
$\Sigma(b)$ is replaced by a bound $\Sh(x_b)$ valid for all $b\le x_b$. When two
witnesses sieve each other's parts, the terms can be re-paired so that each pair is
bounded by twice the maximum over buckets.

\subsection{Revised Lemma C}
Lemma C only uses that $A^*$ lies in a domain $D$ whose classes mod odd $Q$ have
$L_{\rm eff}/Q\pm\mu$ elements. In case II, with $b$ the least even element of $A^*$,
$L_{\rm eff}=N-(b-1)/2$ and $\mu=2$; for an odd witness $b$ and $x\equiv b\pmod4$
we have $xb+1\equiv2\pmod4$, which gives Lemma C$'$ on $(b+4\Z)\cap[b,N]$ with
$L_{\rm eff}=(N-b+1)/4$. For each odd $p_2\ne5$, $\mathrm{non}(p_2)$ is the least of
three valid bounds: truncated inclusion--exclusion without the restriction mod 25,
full inclusion--exclusion as ``all classes minus two'', and truncated
inclusion--exclusion class by class. (Truncated bounds are upper bounds only, so
they cannot be used for a subtracted part.) Lemma C1 is replaced by:

\begin{lemma}[C2$'$]
If $b\le x$, $n_S=\max\{\ell: q_1\cdots q_\ell\le\sqrt{x^2+1},\ q_\ell\le z_2\}$ and
$g(p)=\max(\mathrm{whole}(p)-\mathrm{non}(p),0)$, then
$\sum_{p\in S_b,\,p\le z_2}(\mathrm{whole}(p)-\mathrm{non}(p))$ is at most the sum of
the $n_S$ largest values of $g$. Equivalently, for every $\tau\ge0$ it is at most
$n_S\tau+\sum_p\max(g(p)-\tau,0)$.
\end{lemma}

\subsection{Case analysis and certificate}
The case analysis of Section~\ref{sec:cases} is refined (III.2 is split according
to which of the classes $\{7,32,82\}$, $\{18,68,43\}$ are present, with two or
three witnesses), and every bound is nondecreasing in $N$ for fixed parameters, so
it suffices to check $\mathrm{bound}(N_{hi})<(N_{lo}-6)/25$ on each interval. An
exact certificate (C with GMP) passes on 396 intervals covering
$[6\times10^5,1.6\times10^8]$, with worst ratio $0.9594$; a second implementation in
Python rationals agrees.

\section{Small $N$: a checked colouring}\label{sec:small}

Let $V=\{n:\mu(n^2+1)=0\}$ and let $G$ be the graph on $V$ with an edge $uv$ when
$\mu(uv+1)=0$. Admissible sets are the cliques of $G[V\cap[1,N]]$, so it suffices
to partition $V\cap[1,N]$ into $c_7(N)$ independent sets.

We build such partitions incrementally in $n$: a new class is opened only when
$n\equiv7\pmod{25}$, and every other vertex is placed into an existing class, if
necessary after moving one conflicting element to another class. Every operation is
logged, and an independent program replays the log, recomputing $V$ by a different
sieve and squarefreeness by complete factorisation. Up to $N=4\times10^8$ this gives
$|V|=42{,}063{,}559$ vertices in $16{,}000{,}000=c_7(4\times10^8)$ classes; both
replays pass.

\section{The Lean formalisation}\label{sec:lean}

\subsection{Statement}
The formal statement is
\begin{alltt}\small
def Admissible (A : Finset \(\mathbb{N}\)) : Prop :=
  \(\forall\) a \(\in\) A, \(\forall\) b \(\in\) A, \(\neg\) Squarefree (a * b + 1)
def c7 (N : \(\mathbb{N}\)) : \(\mathbb{N}\) := (N + 18) / 25
def Statement (N : \(\mathbb{N}\)) : Prop :=
  \(\forall\) A : Finset \(\mathbb{N}\), A \(\subseteq\) Finset.Icc 1 N \(\to\) Admissible A \(\to\) A.card \(\le\) c7 N
theorem Erdos848.erdos848 : \(\forall\) N, Statement N
\end{alltt}
together with \lean{erdos848\_max'} (the maximum is exactly $c_7(N)$) and the form
used by the Formal Conjectures project. \lean{\#print axioms} reports only
\lean{propext}, \lean{Classical.choice} and \lean{Quot.sound}.

\subsection{Trust}
Every computation is a Boolean checker evaluated by \lean{decide +kernel}. No
\lean{native\_decide} is used, so neither the compiler nor the interpreter is
trusted. The data (interval parameters, bucket bounds, witness structures, colouring
operations) are produced by untrusted programs; the kernel checks the data, and for
each checker a soundness theorem connects its output to the real bounds of the
lemmas above. As an independent check, \lean{leanchecker}, which is shipped with
Lean, re-checks the declarations of each module of the project in the kernel,
module by module; every module that contains declarations passes. The split in Lean is:
\begin{itemize}[leftmargin=*]
\item $N\le986{,}746$: the colouring, replayed by the kernel;
\item $986{,}747\le N\le1.6\times10^8$: Section~\ref{sec:strong}, 346 intervals,
  worst ratio $0.9306$;
\item $N\ge1.6\times10^8$: Section~\ref{sec:large}, 2913 certificate rows and the tail.
\end{itemize}

\subsection{Kernel arithmetic}
\begin{itemize}[leftmargin=*]
\item \emph{Primes.} For $n<4473^2$, $n$ is prime iff $\gcd(n,\prod_{p\le4472}p)=1$.
  The prime counts and the sums $\sum\lceil10^{30}/p^2\rceil$ at the 10{,}998 points
  needed below $2\times10^7$ are verified by binary-split segment sums. Lookups use
  a search tree written as a literal and checked against the table by one in-order
  traversal.
\item \emph{Squarefreeness} for $n<10007^3$: with $g=\gcd(n,\prod_{p\le10^4}p)$,
  $n$ is squarefree if $\gcd(g,n/g)=1$ and $n/g$ is not a square $>1$.
\item \emph{Lemma C in fixed point.} The sums over the primes $p\le z_2$ are computed
  with scale $2^{40}$; per-prime constants are rounded once, up, or down where they
  are subtracted, and each bucket then needs only natural-number arithmetic. The
  threshold $\tau$ of Lemma C2$'$ is the $n_S$-th largest rounded gain; any
  $\tau\ge0$ is sound. Exact rational sums would involve denominators that are
  products of many prime squares; avoiding them made the evaluation of Lemma C
  about 50 times faster and cut the memory of an interval check by an order of
  magnitude.
\item \emph{The table of Lemma E2} is restricted to the 268 entries with period
  $P\le3000$, and each $\beta_P$ is verified over one period; witness structures
  distinguish odd prime factors $\le23$.
\item \emph{The colouring} for $N\le986{,}746$: an untrusted first-fit proposer with
  augmenting moves emits 117{,}389 operations (open a class, place a vertex, move a
  vertex). The kernel replays them on classes stored in a binary tree of depth 16,
  in 40 windows whose end states are compared with tree-literal checkpoints. The
  vertex set is certified by roots of $x^2\equiv-1\pmod{p^2}$ for every prime
  $p\equiv1\pmod4$, $p\le986{,}746$, together with a proof that these are all such
  primes.
\end{itemize}

\subsection{Cost}
The kernel keeps every reduction of a declaration until the declaration is checked,
so large checks are split into many declarations and reassembled by small lemmas.
A full build takes about five CPU hours (about one hour on a 20-core machine). The
peak memory per module is at most $3.5$~GB for an interval of
Section~\ref{sec:strong}, at most $9.3$~GB for a colouring window, and at most
$7.6$~GB otherwise.

\section{Reproducibility}
All code, certificates and the Lean project are in the accompanying repository,
\url{https://github.com/MacchaCat/erdos848}.
The certificates of Sections~\ref{sec:large} and~\ref{sec:strong} can be re-checked
with the exact Python and C programs; the colouring logs up to $4\times10^8$ are
published with SHA-256 checksums; and \lean{lake build} in the Lean project
re-checks everything in the kernel.

\section*{Disclosure and review status}
The proof, the code and the Lean formalisation were produced mainly by an AI
system (Claude, Anthropic), under the direction of the author, who set the goals and
chose the directions, with a small number of manual interventions. The written
argument went through several rounds of AI review and an independent
re-implementation of all bounds from the text alone; the issues found were fixed.
The complete argument is checked by the Lean kernel. The mathematics has not yet
been reviewed by human experts, and readers should also check that the formal
statement above is the intended one.

This manuscript is released under the Creative Commons Attribution 4.0
International licence (CC BY 4.0); the code and the Lean project are released
under the Apache License 2.0.

\begin{thebibliography}{9}
\bibitem{EP848} T.~F.~Bloom, \emph{Erd\H{o}s Problem \#848},
  \url{https://www.erdosproblems.com/848}.
\bibitem{Sawhney} M.~Sawhney, note on Erd\H{o}s Problem \#848,
  \url{https://www.math.columbia.edu/~msawhney/Problem_848.pdf}.
\bibitem{AI2511} arXiv:2511.16072, Section~IV.1.
\bibitem{Li} A.~C.~Li, \emph{The Exact Extremal Bound in Erd\H{o}s Problem 848},
  August 2026, \url{https://ssrn.com/abstract=7230480}.
\bibitem{LiRepo} A.~C.~Li, \emph{erdos-848-squarefree-product},
  \url{https://github.com/crabsatellite/erdos-848-squarefree-product}.
\bibitem{Lean} L.~de~Moura and S.~Ullrich, \emph{The Lean 4 theorem prover and
  programming language}, CADE-28, LNCS 12699 (2021), 625--635.
\bibitem{Mathlib} The mathlib Community, \emph{The Lean mathematical library},
  CPP 2020, 367--381.
\end{thebibliography}

\end{document}
